%------------------------------------------------------------------------------%
\begin{question}
  Analisando a função \(f(x, y) = x^2 + \frac{y^2}{2}\)
  no ponto \((1, 1)\)

  Encontre as direções nas quais:

  a)\; $f$ cresce mais rapidamente

  b)\; $f$ decresce mais rapidamente

  c)\; $f$ não varia
\end{question}

%------------------------------------------------------------------------------%
\begin{resolution}{Solução}
  Calculando as derivadas parciais

  \[
    \pause
    \D{f}{x}
    \pause
    \:=\; \D{}{x}\left(x^2 + \frac{y^2}{2}\right)
    \pause
    \:=\; 2x
  \]

  \[
    \pause
    \D{f}{y}
    \pause
    \:=\; \D{}{y}\left(x^2 + \frac{y^2}{2}\right)
    \pause
    \:=\; y
  \]
\end{resolution}

%------------------------------------------------------------------------------%
\begin{resolution}{Solução}
  Montando o gradiente
  \pause
  \[
    \nabla f
    \pause
    = \Vetor{\D{f}{x}}{\D{f}{y}}
    \pause
    = \Vetor{\D{}{x}\left(x^2 + \frac{y^2}{2}\right)}
            {\D{}{y}\left(x^2 + \frac{y^2}{2}\right)}
    \pause
    = \vetor{2x}{y}
  \]
  \pause
  Avaliando o gradiente no ponto solicitado
  \[
    \pause
    \nabla f(1, 1)
    \pause
    \;=\; \vetor{2x}{y}\evalat{(1, 1)}{}
    \pause
    \;=\; \vetor{2}{1}
  \]
\end{resolution}

%------------------------------------------------------------------------------%
\begin{resolution}{Solução}
  a)\; $f$ cresce mais rapidamente
  \pause
  na direção do gradiente
  \[
    \pause
    v_M
    \pause
    = \nabla f(1, 1)
    \pause
    = \vetor{2}{1}
  \]

  \pause
  b)\; $f$ decresce mais rapidamente
  \pause
  na direção oposta ao gradiente
  \[
    \pause
    v_m
    \pause
    = -\nabla f (1, 1)
    \pause
    = \vetor{-2}{-1}
  \]
\end{resolution}

%------------------------------------------------------------------------------%
\begin{resolution}{Solução}
  \onslide<+->{c)\; $f$ não varia}
  \onslide<+->{na direção, $v_o$, ortogonal ao gradiente}
  \begin{align*}
    \onslide<+->{\nabla f(1, 1) \cdot v_o            & = 0  \\}
    \onslide<+->{\vetor{2}{1} \cdot \vetor{v_1}{v_2} & = 0  \\}
    \onslide<+->{2v_1 + v_2                          & = 0  \\}
    \onslide<+->{v_2                                 & = -2v_1}
  \end{align*}

  \onslide<+->{As direções são
  \(
    v_o = \vetor{1}{-2}
  \)}
  \onslide<+->{e
  \(
    v_o = \vetor{-1}{2}
  \)}
\end{resolution}

%------------------------------------------------------------------------------%
